\subsection{Υπολογισμός αεροδυναμικών φορτίσεων}

Τα αεροδυναμικά φορτία εξαρτώνται αποκλειστικά από το σχήμα του οχήματος και μπορούν να προσεγγισθούν
με διάφορους τρόπους, ανάλογα με την ακρίβεια που θέλει να πετύχει κανείς και τα φαινόμενα που θέλει να
καταγράψει. Σε αγωνιστικές εφαρμογές, τα αεροδυναμικά φορτία είναι πολύ ευαίσθητα σε αλλαγές του ύψους του
οχήματος από το έδαφος κ.α. Συνήθως, λόγω της πολυπλοκότητάς τους, μοντελοποιείται με πίνακες ή διαγράμματα
μεταβολής που μπορεί να έχουν προκύψει από αναλύσεις ρευστομηχανικής ή πειραματικές μετρήσεις.

Η γενική μορφή της σχέσης του αεροδυναμικού φορτίου δίνεται στην \cref{eq:genericAeroLoad}.

\begin{equation}
   F_{aero} = \dfrac{1}{2}\rho CAU^2
   \label{eq:genericAeroLoad} 
\end{equation}

\noindent
Στην παρούσα διπλωματική εργασία χρησιμοποιείται η μοντελοποίηση των Tremblett \& Limebeer \cite{Tremlett2016} για σκοπούς σύγκρισης. Τα αεροδυναμικά φορτία,
και η κατανομή τους στον εμπρόσθιο και οπίσθιο άξονα μοντελοποιείται με τη γενική
σχέση \ref{eq:genericAeroLoad} και οι εκάστοτε συντελεστές $C$ με ένα πολυώνυμο
δευτέρου βαθμού.

\begin{align}
F_{x,drag} &= -\dfrac{1}{2}C_D\rho AU^2 \label{eq:drag}\\
F_{z,f,lift} &= \dfrac{1}{2}C_L^f\rho AU^2 \label{eq:liftFront}\\
F_{z,r,lift} &= \dfrac{1}{2}C_L^r\rho AU^2
\label{eq:liftRear}
\end{align}

\begin{table}[htbp]
    \centering
    \begin{tabular}{lccc}
        \toprule
        Συντελεστής & $\alpha_0$ & $\alpha_1$ & $\alpha_2$ \\
        \midrule
        $C_L^f$ & $1.614$ & $-1.392 \times 10^{-3}$ & $-4.165 \times 10^{-5}$ \\
        $C_L^r$ & $1.935$ & $\phantom{-}8.335 \times 10^{-3}$ & $-1.249 \times 10^{-4}$ \\
        $C_D$   & $1.053$ & $-6.643 \times 10^{-4}$ & $-1.018 \times 10^{-5}$ \\
        \bottomrule
    \end{tabular}
    \caption{Ενδεικτικοί συντελεστές αεροδυναμικών φορτίων $C = \alpha_2U^2+\alpha_1U+\alpha_0$ \cite{Tremlett2016}.}
    \label{tab:aero_coeffs}
\end{table}